\chapter{变量与类型}
\label{ch:03}

第 1、2 章你已经写过不少变量了，但都是「照抄」。这一章把它讲清楚：

\begin{itemize}

\item 变量是什么、为什么每个变量都要写类型；

\item Magel 有哪些基础类型；

\item 数字、字符、字符串分别怎么写；

\item 以及本章的重头戏：Magel 什么时候替你转类型、什么时候非要你签字。

\end{itemize}

这一章会让你第一次觉得 Magel「有点死板」。请忍着别关页面——第 3.8 节会把这份死板值不值讲清楚。

\section{变量：给值起个名字}

\par
\medskip
\begin{lstlisting}[style=magel]
int64 score = 30;
\end{lstlisting}
\par
\medskip

这一行的读法从左到右：

\begin{center}
\begin{tabular}{ll}
\toprule
部分 & 含义 \\
\midrule
\texttt{int64} & 类型：这个变量装 64 位有符号整数 \\
\texttt{score} & 名字：你自己起的，区分大小写 \\
\texttt{=\allowbreak{}} & 赋值：把右边算出来的值放进这个变量 \\
\texttt{30} & 初始值 \\
\bottomrule
\end{tabular}
\end{center}

类型写在名字前面，这是 C / Java / C++ 的习惯，不是 Rust / TypeScript 那种\texttt{score:\allowbreak{} int64} 或者 \texttt{let score}。和 C 系一样的写法是故意的：Magel 不打算让你重新学一遍已经会的东西。

三个必须记住的细节：

\begin{enumerate}

\item 分号结尾。Magel 不用换行来表示语句结束，写成一整行也行。

\item 区分大小写。\texttt{score} 和 \texttt{Score} 是两个不同的名字。

\item 变量必须先声明再用。写到没声明过的名字是编译错误，不是「默认 0」。

\end{enumerate}

\section{为什么非要写类型}

很多现代语言会帮你推：\texttt{let score =\allowbreak{} 30} 就够了。Magel 能推（下一节就讲 \texttt{auto}），但声明时默认要求你写出来。理由和它的整体性格一致：

\begin{itemize}

\item 读代码的人不用猜。 \texttt{int64 score =\allowbreak{} get\allowbreak{}Score(\allowbreak{})\allowbreak{};} 一眼就知道这里是整数；\texttt{let score =\allowbreak{} get\allowbreak{}Score(\allowbreak{})\allowbreak{}} 得跳去看 \texttt{get\allowbreak{}Score} 的签名。

\item 它不猜你的意思。 记住第 2 章那条报错——\texttt{2026} 不能直接交给要 \texttt{uint64} 的函数。Magel 宁可让你写一个 \texttt{(\allowbreak{}uint64)\allowbreak{}}，也不愿意在某处悄悄换一下。

\item 出错点离原因近。 类型写错，报错就在那一行；不写类型，错可能传到十层外面才炸。

\end{itemize}

\section{整型家族}

\par
\medskip
\begin{lstlisting}[style=magel]
int64 main()
{
    int8 tiny = 127;
    int32 normal = 70000;
    int64 huge = 9000000000;
    uint8 byte = 255;
    println("{} {} {} {}", tiny, normal, huge, byte);
    return 0;
}
\end{lstlisting}
\par
\medskip

\par
\medskip
\begin{lstlisting}[style=out]
127 70000 9000000000 255
\end{lstlisting}
\par
\medskip

\begin{center}
\begin{tabular}{llll}
\toprule
有符号 & 无符号 & 位数 & 大概装得下 \\
\midrule
\texttt{int8} & \texttt{uint8} & 8 & ±127 / 0 到 255 \\
\texttt{int16} & \texttt{uint16} & 16 & ±32767 / 0 到 65535 \\
\texttt{int32} & \texttt{uint32} & 32 & ±21 亿 / 0 到 42 亿 \\
\texttt{int64} & \texttt{uint64} & 64 & ±9.2\ensuremath{\times}10¹⁸ —— 默认选手 \\
\bottomrule
\end{tabular}
\end{center}

日常直接用 \texttt{int64}。 只有两种情况才去挑更窄的：

\begin{itemize}

\item 要跟硬件、文件格式、二进制协议对齐（这时你知道必须是 8 位还是 16 位）；

\item 要省内存（嵌入式板子上，一个 8 位字段和 64 位字段差 8 倍）。

\end{itemize}

为什么没有 \texttt{int}、\texttt{long} 这种名字？ 因为那些名字的含义会随平台变，而这门语言要「一份逻辑从 MCU 跑到云端」——平台上不确定的东西，类型里不能含糊。

\section{浮点}

\par
\medskip
\begin{lstlisting}[style=magel]
int64 main()
{
    float64 ratio = 3.25;
    float32 single =(float32)1.5;
    println("{} {}", ratio, single);
    return 0;
}
\end{lstlisting}
\par
\medskip

\par
\medskip
\begin{lstlisting}[style=out]
3.25 1.5
\end{lstlisting}
\par
\medskip

\texttt{float64}（双精度）是默认选手，\texttt{float32}（单精度）在数值算法、图形、嵌入式里省内存时用。

\texttt{float64} 一栏里那个 \texttt{(\allowbreak{}float32)\allowbreak{}} 是必需的——直接写 \texttt{float32 single =\allowbreak{} 1.\allowbreak{}5;} 会报「不能隐式转换：float64 -> float32」。原因见第 3.8 节：字面量 \texttt{1.\allowbreak{}5} 的类型是\texttt{float64}，往更窄的类型里放就得自己签字。

\begin{tip}{一个容易上当的地方}
\texttt{123456789.\allowbreak{}0} 打印出来是 \texttt{1.\allowbreak{}23e+08}，\texttt{1000000000.\allowbreak{}0} 打印出来是 \texttt{1e+09}。浮点打印会走「短」的那条路，别拿它当精确输出用。要定格式，用 \texttt{format} 的样式参数（第 12 章）。
\end{tip}

\section{\texttt{bool}、\texttt{char}、\texttt{String}}

\par
\medskip
\begin{lstlisting}[style=magel]
int64 main()
{
    bool ready = true;
    bool done = false;
    char initial = 'M';
    String word = "Magel";
    println("{} {} {} {}", ready, done, initial, word);
    return 0;
}
\end{lstlisting}
\par
\medskip

\par
\medskip
\begin{lstlisting}[style=out]
1 0 M Magel
\end{lstlisting}
\par
\medskip

注意第一眼就有点怪的地方：\texttt{bool} 打印出来是 \texttt{1} 和 \texttt{0}，不是 \texttt{true} / \texttt{false}。它不是「字符串化的布尔」，打印时当整数处理了。要比对文本，得自己拼：

\par
\medskip
\begin{lstlisting}[style=magel]
int64 main()
{
    bool ready = true;
    String label;
    if (ready)
    {
        label = "yes";
    }
    else
    {
        label = "no";
    }
    println("{}", label);
    return 0;
}
\end{lstlisting}
\par
\medskip

\par
\medskip
\begin{lstlisting}[style=out]
yes
\end{lstlisting}
\par
\medskip

\begin{tip}{C 程序员注意}
Magel 没有三目运算符。写 \texttt{ready ? "yes" :\allowbreak{} "no"} 是语法错误，老老实实用 \texttt{if /\allowbreak{} else}。可以使用 \texttt{if (\allowbreak{}条件)\allowbreak{} \{ 值一 \} else \{ 值二 \}} 条件表达式；只求值命中的分支。
\end{tip}

\begin{itemize}

\item \texttt{bool} 只有 \texttt{true} 和 \texttt{false} 两个字面量，没有「非零即真」：

\end{itemize}

\par
\medskip
\begin{lstlisting}[style=out]
magel：if 条件必须是 bool，现在给的是 int64（数值条件请显式写成 `(bool)值` 或比较表达式）
\end{lstlisting}
\par
\medskip

硬要拿数字当布尔，得签字写成 \texttt{(\allowbreak{}bool)\allowbreak{}值}——但十有八九你真正想要的是一个比较表达式。

\begin{itemize}

\item \texttt{char} 是单个字符，用单引号：\texttt{'M'}。它和整数不是一回事，\texttt{char c =\allowbreak{} 65;} 会报错，得写 \texttt{(\allowbreak{}char)\allowbreak{}65}。

\item \texttt{String} 是字符串，用双引号，类型是 \texttt{std:\allowbreak{}:\allowbreak{}lang:\allowbreak{}:\allowbreak{}String} 的简写。它属于核心模块 \texttt{lang}，所以不用 import。字符串的内容是第 12 章的主题。

\end{itemize}

\section{常量 \texttt{const} 与推导 \texttt{auto}}

\par
\medskip
\begin{lstlisting}[style=magel]
int64 main()
{
    const int64 LIMIT = 100;
    auto inferred = LIMIT;
    println("{} {}", LIMIT, inferred);
    return 0;
}
\end{lstlisting}
\par
\medskip

\par
\medskip
\begin{lstlisting}[style=out]
100 100
\end{lstlisting}
\par
\medskip

\texttt{const} 表示「这个值定下来就不改了」。给 const 变量再赋值是编译错误。数组长度、缓冲区大小这类东西应该写成 const——不是为了让编译器放心，是为了让下一个读代码的人放心。

\texttt{auto} 表示「类型你替我推」。什么时候用？当右边已经写得很清楚的时候：

\par
\medskip
\begin{lstlisting}[style=magel]
auto info = describe(7);
// 想省一次手写类型
auto picked = range(1, 6).map(square).collect();
// 类型根本写不出来
\end{lstlisting}
\par
\medskip

第二种情况才是 \texttt{auto} 真正的用武之地——后面第 14、17 章会出现那种「类型长得没法看、写出来反而更难读」的场合。

但别滥用。函数内部的小变量用 \texttt{auto} 无所谓，函数签名、结构体字段、全局变量一律写全类型。一句话：类型是给别人看的，不是给编译器看的。

\section{字面量怎么写}

\par
\medskip
\begin{lstlisting}[style=magel]
int64 main()
{
    println("{} {} {}", 0xFF, b1010, 3.14);
    println("{} {}", 2.5e3, 2.5e - 3);
    println("{}", 1e9);
    return 0;
}
\end{lstlisting}
\par
\medskip

\par
\medskip
\begin{lstlisting}[style=out]
255 10 3.14
2.5e+03 0.0025
1e+09
\end{lstlisting}
\par
\medskip

\begin{center}
\begin{tabular}{ll}
\toprule
写法 & 含义 \\
\midrule
\texttt{42} & 十进制整数 \\
\texttt{0x\allowbreak{}FF} & 十六进制（前缀 \texttt{0x}） \\
\texttt{b1010} & 二进制（前缀 \texttt{b}，不是 \texttt{0b}） \\
\texttt{3.\allowbreak{}14} & 浮点，默认 \texttt{float64} \\
\texttt{2.\allowbreak{}5e3} & 科学计数法：2500.0 \\
\texttt{1e9} & 科学计数法：10⁹（没有小数点也算浮点，只要有指数） \\
\texttt{'M'} & 字符（单引号） \\
\texttt{"\textbackslash\{\}n"} & 字符串（双引号），里面有转义 \\
\bottomrule
\end{tabular}
\end{center}

\begin{tip}{提示}
二进制的前缀是 \texttt{b} 而不是 \texttt{0b}：\texttt{b1010} 是 10，写 \texttt{0b1010} 会报词法错误。这条容易记混，见本章最后的表。\par 带指数的就是浮点。 \texttt{1e9} 和 \texttt{2.\allowbreak{}5e3} 一样是 \texttt{float64}，所以 \texttt{int64 x =\allowbreak{} 1e9;} 会报「不能隐式转换：float64 -> int64」——浮点往整数里放要签字（见 3.9 节）。
\end{tip}

\subsection{数字后缀：\texttt{1u} 与 \texttt{1.\allowbreak{}0f}}

小写 \texttt{u} 表示无符号整数；没有类型上下文时默认是 \texttt{uint64}。无符号目标类型可以按自己的位宽接收这个字面量，并检查取值范围。它不会隐式变成有符号整数，也不能对无符号字面量直接写一元负号。

小写 \texttt{f} 把浮点字面量定为 \texttt{float32}。不带后缀的 \texttt{1.\allowbreak{}0} 是 \texttt{float64}；\texttt{1f}、\texttt{1.\allowbreak{}0f}、\texttt{1e2f} 都是 float32 写法。后缀必须小写。

\par
\medskip
\begin{lstlisting}[style=magel]
int64 floatingKind(float32 value)
{
    return 32;
}

int64 floatingKind(float64 value)
{
    return 64;
}

uint64 unsignedValue(uint64 value)
{
    return value;
}

int64 main()
{
    auto one = 1u;
    auto single = 1.0f;
    uint8 narrow = 255u;
    uint64 maximum = 18446744073709551615u;
    println("{} {} {} {}", unsignedValue(one), floatingKind(single),
            floatingKind(1.0), narrow);
    println("{}", maximum);
    return 0;
}
\end{lstlisting}
\par
\medskip

\par
\medskip
\begin{lstlisting}[style=out]
1 32 64 255
18446744073709551615
\end{lstlisting}
\par
\medskip

\texttt{uint8 narrow =\allowbreak{} 256u;} 会因为超范围报错。\texttt{int64 signed\allowbreak{}Value =\allowbreak{} 1u;} 不能隐式转换，\texttt{-1u} 也会被拒绝；需要有符号负数时直接写 \texttt{-1}。\texttt{0xffu}、\texttt{b\_101u} 同样支持无符号后缀。十六进制和现有二进制写法支持数字间的下划线分隔，但十进制 \texttt{1\_000} 和 \texttt{0b1010} 并不是当前文法支持的写法。

\section{隐式转换：只放行「变宽」}

这是本章最重要的一节。

先说结论，别被别的语言带偏：Magel 只放行一类转换——「变宽」，也就是「目标装得下源、不会悄悄丢东西」的方向。反过来（变窄、有符号放进无符号、浮点放进整数、\texttt{char}/\texttt{bool} 跟整数互放）一律要你签字。所以准确的说法不是「没有隐式转换」，而是只放行「装得下」的那半边。

\par
\medskip
\begin{lstlisting}[style=magel]
int64 main()
{
    int32 a = 100;
    int64 b = a;
    // 整型同符号变宽，放行
    uint8 u = 200;
    int32 m = u;
    // 无符号 -> 更宽的有符号，放行
    float64 c = b;
    // 整型 -> 浮点，放行
    println("{} {} {}", b, m, c);
    return 0;
}
\end{lstlisting}
\par
\medskip

\par
\medskip
\begin{lstlisting}[style=out]
100 200 100
\end{lstlisting}
\par
\medskip

放行清单（不用签字）：

{\small\setlength{\tabcolsep}{4pt}
\begin{longtable}{>{\raggedright\arraybackslash}p{\dimexpr 0.269\linewidth-2\tabcolsep\relax}>{\raggedright\arraybackslash}p{\dimexpr 0.711\linewidth-2\tabcolsep\relax}}
\toprule
方向 & 例子 \\
\midrule
\endhead
同符号整型变宽 & \texttt{int8 → int16/\allowbreak{}int32/\allowbreak{}int64}、\texttt{int32 → int64}、\texttt{uint8 → uint32} …… \\
无符号 \ensuremath{\rightarrow} 更宽的有符号 & \texttt{uint8 → int16/\allowbreak{}int32/\allowbreak{}int64}、\texttt{uint16 → int32}、\texttt{uint32 → int64} \\
整型 \ensuremath{\rightarrow} 浮点 & \texttt{int64 → float64}、\texttt{int32 → float32} \\
浮点变宽 & \texttt{float32 → float64} \\
\bottomrule
\end{longtable}
}

下面这三类，方向是反的，必须自己签字：

\par
\medskip
\begin{lstlisting}[style=magel]
int64 main()
{
    int64 big = 70000;
    int32 small = big;
    // 缩窄：位宽变小
    println("{}", small);
    return 0;
}
\end{lstlisting}
\par
\medskip

\par
\medskip
\begin{lstlisting}[style=out]
magel：不能隐式转换：int64 -> int32（请显式强转）
\end{lstlisting}
\par
\medskip

\begin{itemize}

\item 缩窄：位宽变小（\texttt{int64 → int32}、\texttt{float64 → float32}、\texttt{uint8 → int8}）——装不下，要签字；

\item 有符号 \ensuremath{\rightarrow} 无符号：哪怕目标更宽也不行（\texttt{int32 → uint32}、\texttt{int64 → uint64}）——负数塞不进无符号，要签字；

\item 跨大类：浮点 \ensuremath{\rightarrow} 整型、\texttt{char} / \texttt{bool} \ensuremath{\leftrightarrow} 整型——含义会变，要签字。

\end{itemize}

\begin{tip}{提示}
一句话记法：装得下的方向，Magel 自动帮你；装不下、或者换符号/大类，必须你亲自签字。
\end{tip}

这么一管，麻烦是麻烦了点，但它买到的是三件事：

\begin{enumerate}

\item 每一处精度损失都是有意的。 \texttt{int64} 里装着 90 亿，你要往 \texttt{int32} 里塞——这一下你必然写了 \texttt{(\allowbreak{}int32)\allowbreak{}}，将来出了问题，你能顺着这个强转找到是谁决定的。

\item 重载能被准确选中。 第 2 章那条「找不到匹配的重载 \texttt{(\allowbreak{}int64)\allowbreak{}}」就是这套规则的副产品：因为不替你「凑合着转」，编译器只能说「你要 \texttt{uint64}，给我的是 \texttt{int64}」——它把真相告诉你了，而不是塞进某个它认为差不多的重载。

\item 跨平台的结果是确定的。 在 32 位板子上 \texttt{long} 是 4 字节、PC 上是 8 字节，这种「同一个表达式换台机器结果不同」的坑，Magel 用「类型必须对上」堵掉了。

\end{enumerate}

什么时候你会想骂它？ 写累加、写混合计算的时候。这时候记住两句话：

\begin{itemize}

\item 拿不准就用 \texttt{int64} / \texttt{float64}，让变量类型统一，转换就少了；

\item 真要转，把它写在最显眼的地方（用括号包住，别藏在表达式中间）。

\end{itemize}

\section{显式强转：\texttt{(\allowbreak{}目标类型)\allowbreak{}值}}

\par
\medskip
\begin{lstlisting}[style=magel]
int64 main()
{
    int64 big = 70000;
    int32 normal =(int32)big;
    uint8 byte =(uint8)300;
    char letter =(char)65;
    println("{} {} {}", normal, byte, letter);
    return 0;
}
\end{lstlisting}
\par
\medskip

\par
\medskip
\begin{lstlisting}[style=out]
70000 44 A
\end{lstlisting}
\par
\medskip

写法就是 C 的老规矩：在值前面用括号写上目标类型。

三行分别演示了三件事：

\begin{itemize}

\item \texttt{(\allowbreak{}int32)\allowbreak{}big} —— 缩窄。70000 装得进 \texttt{int32}，值不变。

\item \texttt{(\allowbreak{}uint8)\allowbreak{}300} —— 截断。300 装不进 8 位，只留下低 8 位，得到 \texttt{44}。强转不会帮你检查装得下装不下——签了字就由你负责。

\item \texttt{(\allowbreak{}char)\allowbreak{}65} —— 数字变字符，65 就是 \texttt{'A'}。

\end{itemize}

\begin{tip}{提示}
强转是「你签字、编译器照做」，不做范围检查。 所以 \texttt{(\allowbreak{}int8)\allowbreak{}200} 得到 \texttt{-56}、\texttt{(\allowbreak{}uint8)\allowbreak{}300} 得到 \texttt{44}，都不报错。会做范围检查的是字面量直接落进槽位那条路（3.10 节）——两条路泾渭分明：没写强转，编译器替你把关；写了强转，责任在你。
\end{tip}

反过来 \texttt{(\allowbreak{}int64)\allowbreak{}3.\allowbreak{}9} 得到 \texttt{1}：浮点转整数是直接砍掉小数位，不是四舍五入。要四舍五入得自己想办法（\texttt{std:\allowbreak{}:\allowbreak{}math} 里有取整函数，第 12 章提）。

\subsection{\texttt{cast<\allowbreak{}T>\allowbreak{}(\allowbreak{}value)\allowbreak{}}：把目标类型写清楚}

除了前面的 C 风格强转，也可以写 \texttt{cast<\allowbreak{}目标类型>\allowbreak{}(\allowbreak{}值)\allowbreak{}}。这是显式转换，不是允许普通赋值偷偷缩窄。

\par
\medskip
\begin{lstlisting}[style=magel]
int64 main()
{
    int64 original = 42;
    int32 narrowed = cast<int32>(original);
    println("{}", narrowed);
    return 0;
}
\end{lstlisting}
\par
\medskip

\par
\medskip
\begin{lstlisting}[style=out]
42
\end{lstlisting}
\par
\medskip

尖括号中是目标类型，圆括号中是被转换的表达式。自定义类型的显式转换要提供对应协议；宿主类型的 \texttt{impl} 绑定见第 21 章。转换也不会让一个引用越过来源的生命周期。

\section{字面量的取值范围}

这一节两条规矩，都是本书撰写时实测确认过的：

规矩一：负数字面量可以直接写进窄类型。

\texttt{-300} 在语法上是「一元负号 + 正数字面量 \texttt{300}」，但编译器把它当常量看值——只要值落在目标类型的范围内，就直接放行：

\par
\medskip
\begin{lstlisting}[style=magel]
int64 main()
{
    int16 small = -300;
    // -300 落在 int16（-32768..32767）里
    int8 tiny = -100;
    println("{} {}", small, tiny);
    return 0;
}
\end{lstlisting}
\par
\medskip

\par
\medskip
\begin{lstlisting}[style=out]
-300 -100
\end{lstlisting}
\par
\medskip

规矩二：十进制字面量超出范围会报错，不静默回绕。

\par
\medskip
\begin{lstlisting}[style=magel]
int64 main()
{
    int8 overflow = 200;
    // int8 只到 127
    return 0;
}
\end{lstlisting}
\par
\medskip

\par
\medskip
\begin{lstlisting}[style=out]
magel：整数字面量 200 超出 int8 的取值范围（-128..127）
\end{lstlisting}
\par
\medskip

\texttt{uint8 b =\allowbreak{} 255;} 没问题（255 在 \texttt{0.\allowbreak{}.\allowbreak{}255} 里），\texttt{uint8 b =\allowbreak{} 256;} 报同样的错。这条只管「字面量落进槽位」（初始化 / 赋值）；显式强转是另一回事——\texttt{(\allowbreak{}int8)\allowbreak{}200} 是你亲口签字，编译器只管按位截断（见 3.9 节）。

\begin{tip}{提示}
十六进制 / 二进制字面量按「位模式」处理，不做范围检查。 \texttt{int8 x =\allowbreak{} 0x\allowbreak{}FF;} 得到 \texttt{-1}（低 8 位全 1，按有符号读就是 -1），和 C 一样。要写明确的位模式，就用 \texttt{0x} / \texttt{b} 前缀。
\end{tip}

还有几个词法上的「不认」，撞上了照右列改：

\begin{center}
\begin{tabular}{lll}
\toprule
你可能想写 & 实际会怎样 & 改成 \\
\midrule
\texttt{0b1010} & 词法错误（前缀是 \texttt{b}，不是 \texttt{0b}） & \texttt{b1010} \\
\texttt{1\_000} & 语法错误（下划线分隔不认） & \texttt{1000} \\
\texttt{1.\allowbreak{}0F} / \texttt{1U} & 后缀必须小写 & \texttt{1.\allowbreak{}0f} / \texttt{1u} \\
\bottomrule
\end{tabular}
\end{center}

\section{小结}

{\small\setlength{\tabcolsep}{4pt}
\begin{longtable}{>{\raggedright\arraybackslash}p{\dimexpr 0.078\linewidth-2\tabcolsep\relax}>{\raggedright\arraybackslash}p{\dimexpr 0.902\linewidth-2\tabcolsep\relax}}
\toprule
你会了 & 关键点 \\
\midrule
\endhead
声明 & \texttt{int64 score =\allowbreak{} 30;}，类型在前、名字在后、分号结尾 \\
整型 & \texttt{int8/\allowbreak{}16/\allowbreak{}32/\allowbreak{}64} + \texttt{uint8/\allowbreak{}16/\allowbreak{}32/\allowbreak{}64}，日常用 \texttt{int64} \\
浮点 & 无后缀为 \texttt{float64}，\texttt{1.\allowbreak{}0f} 为 \texttt{float32}，也可显式转换 \\
其它标量 & \texttt{bool} 只有 \texttt{true/\allowbreak{}false}（打印成 1/0）、\texttt{char} 用单引号 \\
字符串 & \texttt{String} 用双引号，属核心模块 \texttt{lang}，不用 import \\
常量 & \texttt{const int64 LIMIT =\allowbreak{} 100;} \\
推导 & \texttt{auto}，适合右边已经把类型写清楚的场合 \\
字面量 & \texttt{0x\allowbreak{}FF}、\texttt{b1010}、\texttt{3.\allowbreak{}14}、\texttt{2.\allowbreak{}5e3}、\texttt{1e9}、\texttt{'M'}、\texttt{"\textbackslash\{\}n"} \\
转换 & 只放行装得下的方向（同符号整型变宽、无符号\ensuremath{\rightarrow}更宽的有符号、整\ensuremath{\rightarrow}浮、\texttt{float32}\ensuremath{\rightarrow}\texttt{float64}）；其余要显式强转 \texttt{(\allowbreak{}int64)\allowbreak{}x} \\
字面量范围 & 十进制字面量超范围直接报错（\texttt{int8 x =\allowbreak{} 200;} 不合法）；\texttt{0x}/\texttt{b} 字面量按位模式放行 \\
\bottomrule
\end{longtable}
}

一句话总结这一章：Magel 要把「类型在哪里变了」留成一件看得见的事。

\section{下一步}

现在你会声明变量、会算数，但所有代码都堆在 \texttt{main} 里。第 2 章那个猜数字程序要是再长一点，\texttt{main} 就没法看了。

下一章解决这个问题：把代码切成函数。

翻到\hyperref[ch:04]{第 4 章 函数}。
